\documentclass[tikz,border=6pt]{standalone}
\usepackage{ctex}
\usepackage{amsmath}
\usepackage{amssymb}
\usepackage{pgfplots}
\pgfplotsset{compat=1.18}
\usepgfplotslibrary{fillbetween}
\usetikzlibrary{calc,angles,quotes,arrows.meta,positioning,decorations.markings}
\begin{document}
\begin{tikzpicture}
\begin{axis}[
    width=12cm, height=9cm,
    xlabel={$x$}, ylabel={$y$},
    xmin=-2.4, xmax=1.8, ymin=-0.6, ymax=6,
    xtick={-2,-1,0,1},
    ytick={0,1,2,3,4,5,6},
    axis lines=middle,
    grid=major, grid style={gray!18},
    tick label style={font=\small},
    label style={font=\small},
    legend style={font=\small, draw=gray!40, at={(0.03,0.97)}, anchor=north west},
    enlarge x limits=0.02,
]
% y = e^x 曲线
\addplot[blue!75!black, very thick, name path=curve, domain=-2.4:1.8, samples=120] {exp(x)};
\addlegendentry{$y=e^x$（严格凸）}
% 切线 y = 1+x （x=0 处）
\addplot[red!80!black, very thick, name path=tang, domain=-2.4:1.8, samples=2] {1+x};
\addlegendentry{切线 $y=1+x$（$x=0$ 处）}
% 曲线与切线之间的间隙填充：凸性使曲线恒在切线上方
\addplot[gray!25, opacity=0.7] fill between[of=curve and tang];
% 切点
\addplot[only marks, mark=*, mark size=2pt, black] coordinates {(0,1)};
\node[anchor=south east, font=\small] at (axis cs:0,1) {切点 $(0,1)$};
% 标注：切线恒在曲线下方
\node[anchor=west, align=left, font=\small,
      fill=white, fill opacity=0.8, text opacity=1, inner sep=3pt, draw=gray!30]
    at (axis cs:0.18,3.4)
    {切线恒在曲线下方\\$\Rightarrow\ e^x\ge 1+x$\\（仅 $x=0$ 取等）};
% 在 x=-1.5 处用竖线示意 gap（切线已落到负值，曲线仍为正）
\addplot[gray!55, dashed, thick] coordinates {(-1.5,-0.5) (-1.5,0.2231)};
\node[anchor=north, font=\footnotesize, red!70!black] at (axis cs:-2.0,-0.32) {切线值 $1+x<0$};
\node[anchor=south, font=\footnotesize, blue!60!black]
    at (axis cs:-1.5,0.30) {$e^{-1.5}\approx0.22>0$};
\end{axis}
\end{tikzpicture}
\end{document}
